\section{Results}
\label{sec:results}

% ------------------------------------------------------------------
\subsection{Main Comparison}
\label{sec:results:main}

\Cref{tab:main_results} presents the final evaluation returns for all methods.
\todo{Fill in with actual numbers.}

\begin{table}[t]
  \centering
  \caption{Mean evaluation return ($\pm$ std) over 3 seeds after $1\text{M}$
           environment steps.  Best result per environment in \textbf{bold}.}
  \label{tab:main_results}
  \begin{tabular}{@{}lcc@{}}
    \toprule
    \textbf{Method} & \textbf{HalfCheetah-v5} & \textbf{Hopper-v5} \\
    \midrule
    Holistic PPO              & \todo{XXX $\pm$ YYY} & \todo{XXX $\pm$ YYY} \\
    Per-Step Uniform          & \todo{XXX $\pm$ YYY} & \todo{XXX $\pm$ YYY} \\
    Per-Step Learned Global   & \todo{XXX $\pm$ YYY} & \todo{XXX $\pm$ YYY} \\
    Per-Step State-Dependent  & \todo{XXX $\pm$ YYY} & \todo{XXX $\pm$ YYY} \\
    \bottomrule
  \end{tabular}
\end{table}

\todo{Discuss which methods outperform the baseline and by how much.  Note any
environment-specific differences.}

% ------------------------------------------------------------------
\subsection{Training Curves}
\label{sec:results:curves}

\Cref{fig:training_curves} shows the evaluation return over the course of
training.

\begin{figure}[t]
  \centering
  \begin{subfigure}[t]{0.48\textwidth}
    \centering
    \includegraphics[width=\textwidth]{figs/halfcheetah_training_curves.pdf}
    \caption{HalfCheetah-v5}
    \label{fig:curve_halfcheetah}
  \end{subfigure}
  \hfill
  \begin{subfigure}[t]{0.48\textwidth}
    \centering
    \includegraphics[width=\textwidth]{figs/hopper_training_curves.pdf}
    \caption{Hopper-v5}
    \label{fig:curve_hopper}
  \end{subfigure}
  \caption{Evaluation return vs.\ environment steps for all methods.  Solid
           lines show the mean over 3 seeds; shaded regions show $\pm 1$ std.}
  \label{fig:training_curves}
\end{figure}

\todo{Describe learning speed differences.  Does per-step PPO learn faster or
more stably?}

% ------------------------------------------------------------------
\subsection{Weight Profile Analysis}
\label{sec:results:weights}

For the learned-global weighting scheme, \Cref{fig:weight_profile} visualizes
the learned weight distribution $w_k$ across the $K$ flow steps after
training.

\begin{figure}[t]
  \centering
  \includegraphics[width=0.5\textwidth]{figs/weight_profile.pdf}
  \caption{Learned global weights $w_k$ (softmax of $\boldsymbol{\alpha}$)
           across flow steps for HalfCheetah-v5.  \todo{Describe pattern.}}
  \label{fig:weight_profile}
\end{figure}

\todo{Discuss which steps receive the most credit.  Do early steps (close to
the final action) or late steps (close to noise) get higher weight?  Is the
pattern consistent across environments?}

% ------------------------------------------------------------------
\subsection{Per-Step Clip Fraction Analysis}
\label{sec:results:clipfrac}

\Cref{fig:clip_fractions} compares the per-step clip fractions (fraction of
samples where $r_{t,k}$ is clipped) between holistic and per-step PPO.

\begin{figure}[t]
  \centering
  \includegraphics[width=0.5\textwidth]{figs/clip_fractions.pdf}
  \caption{Per-step clip fractions during training.  \todo{Describe the
           difference between holistic and per-step clipping.}}
  \label{fig:clip_fractions}
\end{figure}

\todo{Discuss whether per-step clipping leads to more uniform and moderate clip
fractions compared to holistic clipping.}

% ------------------------------------------------------------------
\subsection{Ablation: Number of Flow Steps $K$}
\label{sec:results:ablation_K}

\Cref{tab:ablation_K} shows results as we vary the number of flow steps.

\begin{table}[t]
  \centering
  \caption{Evaluation return on HalfCheetah-v5 for different $K$.}
  \label{tab:ablation_K}
  \begin{tabular}{@{}lccc@{}}
    \toprule
    \textbf{Method} & $K{=}2$ & $K{=}4$ & $K{=}8$ \\
    \midrule
    Holistic PPO              & \todo{---} & \todo{---} & \todo{---} \\
    Per-Step Uniform          & \todo{---} & \todo{---} & \todo{---} \\
    Per-Step Learned Global   & \todo{---} & \todo{---} & \todo{---} \\
    \bottomrule
  \end{tabular}
\end{table}

\todo{Discuss how the gap between holistic and per-step changes with $K$.
Hypothesis: per-step PPO should benefit more for larger $K$, where the holistic
ratio becomes more extreme.}

% ------------------------------------------------------------------
\subsection{Ablation: Noise Scale $\sigma$}
\label{sec:results:ablation_sigma}

\begin{table}[t]
  \centering
  \caption{Evaluation return on HalfCheetah-v5 for different noise scales
           $\sigma$, using per-step uniform PPO with $K{=}4$.}
  \label{tab:ablation_sigma}
  \begin{tabular}{@{}lccc@{}}
    \toprule
    \textbf{Method} & $\sigma{=}0.01$ & $\sigma{=}0.1$ & $\sigma{=}0.5$ \\
    \midrule
    Holistic PPO     & \todo{---} & \todo{---} & \todo{---} \\
    Per-Step Uniform & \todo{---} & \todo{---} & \todo{---} \\
    \bottomrule
  \end{tabular}
\end{table}

\todo{Discuss sensitivity to $\sigma$.  Smaller $\sigma$ concentrates density
and may produce sharper ratios; larger $\sigma$ makes per-step densities
broader.}

% ------------------------------------------------------------------
\subsection{Discussion}
\label{sec:results:discussion}

\todo{Synthesize findings.  When does per-step PPO help most?  Is
state-dependent weighting worth the additional complexity?  What are failure
modes?}
